\section{The Demons' Powers and Their Limits}\label{app:demons}

The demons keep the natural powers of the angelic nature (\ST{I q.~64
a.~1}); their will is fixed in evil (a.~2); and every use of their power
falls under the command, permission, or restraint of God. Gregory states the
rule: ``the will of Satan is always evil, but his power is never unjust, for
his will he derives from himself, but his power he derives from God''
(\emph{Moralia} II.10.17). Augustine gives the double limit, of permission
and of nature: a man ``can walk, yet~\dots\ cannot do so if he is not
permitted; but~\dots\ cannot fly, even if he be permitted'' (\emph{De Trin.}
III.9.18). The rites and discipline of exorcism belong to the study \emph{Catholic Exorcism: History, Discipline, and
Pastoral Practice}.

\begingroup
\footnotesize
\setlength{\tabcolsep}{3pt}
\begin{longtable}{@{}>{\bfseries\raggedright\arraybackslash}p{0.12\linewidth}>{\raggedright\arraybackslash}p{0.22\linewidth}>{\raggedright\arraybackslash}p{0.23\linewidth}>{\raggedright\arraybackslash}p{0.13\linewidth}>{\raggedright\arraybackslash}p{0.22\linewidth}@{}}
\toprule
\textbf{Power} & \textbf{What the demons can do} & \textbf{Limit under God} &
\textbf{Aquinas} & \textbf{Fathers and the Church}\\
\midrule
\endhead
\bottomrule
\endfoot
Natural knowledge & Keep their natural knowledge whole; the angelic gifts
``exist, and are complete, and all luminous'' in them & Speculative
knowledge by grace is lessened to what is revealed; affective knowledge and
charity are lost; they knew the Incarnation less fully; error is possible in
what exceeds nature & \ST{I q.~64 a.~1}; \emph{De malo} q.~16 a.~6 &
Dionysius, \emph{DN} 4.23; Damascene, \emph{De fide orth.} II.4 (``they
know the Scriptures''); Augustine, \emph{De div.\ daem.} 6.10\\
Knowledge of the future & Foreknow by conjecture what proceeds from causes
usually or necessarily; foretell by speed and observation, and often what
they themselves are about to do & The future in itself is known to God
alone; they are deceived when a higher command overturns their plans, and
they deceive from envy & \ST{I q.~57 a.~3}; \emph{De malo} q.~16 a.~7 &
Athanasius, \emph{Vita Antonii} 31--33; Damascene, \emph{De fide orth.}
II.4; Augustine, \emph{De div.\ daem.} 5.9--6.10\\
Knowledge of thoughts & Learn thoughts in their effects, from signs in the
body and bearing, and explore a man's disposition by temptation & The
thought as it is in the will is known only to God and to the one who wills
& \ST{I q.~57 a.~4}; q.~114 a.~2 ad~2; \emph{De malo} q.~16 a.~8 &
Cassian, \emph{Conl.} VII.15--16; Augustine, \emph{De div.\ daem.} 5.9,
qualified in \emph{Retract.} II.30\\
Imagination and senses & Move the imagination through the spirits and
humours, sleeping or waking; present sensible objects outwardly; change the
senses from within; rouse the passions & Cannot impress a form never
received from sense (``a man born blind''); cannot infuse species into the
intellect; the will remains free to consent or resist & \ST{I q.~111
aa.~2--4}; \emph{De malo} q.~16 aa.~11--12 & Augustine, \emph{De civ.\ Dei}
XVIII.18; Cassian, \emph{Conl.} VII.8; Origen, \emph{De princ.} III.2.4\\
Local motion of bodies & Move bodies locally, including bodies not joined to
them, and through local motion apply natural agents & Cannot induce a
substantial form immediately; create nothing; do not rule the weather by
their own authority; act only as far as permitted & \ST{I q.~110 aa.~2--3};
\emph{De malo} q.~16 aa.~9--10 & Braga~I, can.~8 (DS~458); Augustine,
\emph{De Trin.} III.8.13, III.9.18; Damascene, \emph{De fide orth.} II.3--4\\
Wonders & Produce real effects by applying the ``seeds'' in the elements, as
the serpents and frogs of Pharaoh's magicians and the fire and wind against
Job; work illusions on the senses & No miracle in the proper sense; no man
turned into a beast and no dead man raised; the magicians failed at the
gnats, ``the finger of God'' & \ST{I q.~110 a.~4}; q.~114 a.~4; \emph{SCG}
III.103; \emph{De pot.} q.~6 a.~5 & Augustine, \emph{De Trin.}
III.7.12--8.13, \emph{De civ.\ Dei} XVIII.18, XX.19; Exod 8:18--19\\
Temptation & Tempt to sin as their proper office; persuade and kindle
thoughts; the devil is the indirect cause of all sins through the first sin
& Cannot force the will; not the direct cause of every sin; not beyond a
man's strength; the Powers restrain them; the conquered withdraw for a time
& \ST{I q.~111 a.~2}; q.~114 aa.~1--3, 5 & Origen, \emph{De princ.}
III.2.2--3; Cassian, \emph{Conl.} VII.8, 20, 22; Gregory, \emph{Moralia}
II.10--11, \emph{Hom.\ in Ev.} 34.10; 1~Cor 10:13\\
Dominion over fallen man & By Adam's sin man fell into captivity under the
devil's power & Broken by Christ's victory; Satan's power is a creature's
and not infinite, and permitted by providence & \ST{I q.~114 a.~1} &
Trent, DS~1511, 1521; CCC~395, 2853; Athanasius, \emph{Vita Antonii} 28,
41; Col 2:15\\
\end{longtable}
\endgroup

The remedies the Fathers name are the weapons of the Christian life.
Antony sums them up: the demons ``fear the fasting, the
sleeplessness, the prayers, the meekness, the quietness, the contempt of
money and vainglory, the humility, the love of the poor, the alms, the
freedom from anger of the ascetics, and, chief of all, their piety towards
Christ'' (\emph{Vita Antonii} 30). Paul~VI gave the same answer:
\latin{tutto ciò che ci difende dal peccato ci ripara per ciò stesso
dall'invisibile nemico. La grazia è la difesa decisiva}, whatever defends us
from sin shields us by that very fact from the invisible enemy; grace is the
decisive defence (audience of 15 November 1972).

\begin{fourstable}{Remedy}{Scripture (Douay)}{Fathers}{Aquinas and the
Church}
The Cross & ``despoiling the principalities and powers'' (Col 2:15) &
Athanasius, \emph{Vita Antonii} 23, 35 (``they are cowards, and greatly fear
the sign of the Lord's Cross''); \emph{De incarnatione} 47.2 (``by the Sign
of the Cross, though a man but use it, he drives out their deceits'');
Cassian, \emph{Conl.} VII.23, VIII.18 & CCC~2853: the victory ``won once for
all at the Hour'' of the Passion\\
The name of Jesus & ``The devils also are subject to us in thy name'' (Luke
10:17); ``In my name they shall cast out devils'' (Mark 16:17); Phil 2:10 &
Athanasius, \emph{Vita Antonii} 40--41; Origen, \emph{Contra Celsum} I.6
(``Such power, indeed, does the name of Jesus possess over evil
spirits''); Augustine, \emph{De div.\ qq.\ 83} q.~79.4 & \ST{I q.~110 a.~4
ad~2}: the signs of public justice, ``as by invoking the name of Christ''\\
Prayer and fasting & ``This kind can go out by nothing, but by prayer and
fasting'' (Mark 9:28[29]); ``Watch ye: and pray that ye enter not into
temptation'' (Matt 26:41); Matt 6:13 & Athanasius, \emph{Vita Antonii} 23
(``by prayer, fasting, and faith in the Lord their attack immediately
fails''), 30 & CCC~2849: the battle and the victory ``become possible only
through prayer''; Paul~VI, 15 November 1972, on Mark 9:28[29]\\
Faith, humility, and resistance & ``Taking the shield of faith'' (Eph
6:16); ``Whom resist ye, strong in faith'' (1~Pet 5:9); ``resist the devil:
and he will fly from you'' (James 4:7) & Antony's counsel to rejoice in the
Lord and despise the enemy (\emph{Vita Antonii} 42); Cassian, \emph{Conl.}
VII.8; Origen, \emph{De princ.} III.2.3--4 & \ST{I q.~111 a.~2}: the will
``ever remains free to consent to, or to resist''\\
Grace and the sacraments & ``that, through death, he might destroy him who
had the empire of death'' (Heb 2:14) & Chrysostom, \emph{Hom.\ in Ioann.}
46: ``Let us then return from that table like lions breathing fire, having
become terrible to the devil''; Ambrose to the newly baptized, in CCC~2852
& Trent on the captivity from which justification frees (DS~1511, 1521);
\ST{I q.~114 a.~1 ad~2}: the recompense gained ``principally through the
grace of God''\\
The holy angels & ``Fear not: for there are more with us than with them''
(4~Kgs 6:16); Ps 33[34]:8; Ps 90[91]:11 & Origen, \emph{De princ.} III.2.5
(the angel who stands with Jacob in the struggle); Gregory, \emph{Hom.\ in
Ev.} 34.10 (the Powers bridle the adverse powers), \emph{Moralia}
XVII.13.19 & \ST{I q.~114 a.~1 ad~2}: ``secondarily through the
guardianship of the angels''; \S\ref{sec:guardians}\\
\end{fourstable}
